\begin{helper}[Terminal component after the outside swap]
In a defective-color setup, fix a defective color $d$, put
$\gamma=\gamma( d)$, and interchange $\alpha_0$ and $\gamma$ on every
edge not incident with $V(T)$, obtaining $\varphi_1$.  Then
$\varphi_1$ is proper, has exactly the original uncolored identity, is
stable on $T$, and fixes every edge incident with $V(T)$.

Let $K$ be the $(\gamma,d)$-component of $v^0_d$ under $\varphi_1$.
Its support is exactly the set of nodes of the recorded path $P_d$ at or
after the recorded occurrence of $v^0_d$.  In particular $f^0_d\in K$,
\[
 K\cap\{\operatorname{inl}v:v\in V(T)\}
   =\{\operatorname{inl}v^0_d\},
\]
no displayed edge identity of $T$ belongs to $K$, and every edge-node in
$K$ has color $\gamma$ or $d$ under $\varphi_1$.
\end{helper}
\begin{proof}
We first expose the hypotheses hidden in the configuration records.  By
\noderef{ScheideReducedConfiguration}, $C$ contains an inherited
Scheide structural configuration.  The latter definition,
\noderef{ScheideStructuralConfiguration}, supplies a globally maximum
triple.  Unpacking \noderef{GlobalMaximumTriple} supplies in turn a
$k$-triple, and \noderef{KTriple} says in particular that $\varphi$ is
proper in the sense of \noderef{IsProperEdgeColoring} and leaves exactly
$e_0$ uncolored in the sense of \noderef{IsOneEdgeUncolored}.  Applying
\noderef{ScheideGlobalMaximumOrder} to that globally maximum triple shows
that $V(T)$ is closed under $\varphi$.

Put $y_0=C.y(0)$ and $v_0=v^0_d$.  The data in
\noderef{ScheideDefectiveSetup} say that $\alpha_0$ is missing at $y_0$,
that $\gamma$ is missing at $v_0$, that both vertices lie in $V(T)$, and
that $\alpha_0\ne\gamma$.  Hence both colors belong to the union
\noderef{TreeMissingColors}.  We have now checked every hypothesis of
\noderef{OutsideTreeColorSwapStable}: properness and the unique uncolored
identity come from the displayed chain through \noderef{KTriple},
closedness comes from \noderef{ScheideGlobalMaximumOrder}, and missingness
and distinctness come from \noderef{ScheideDefectiveSetup}.  That lemma
therefore proves that $\varphi_1$ is proper, leaves exactly $e_0$
uncolored, is $(T,\varnothing,\varphi)$-stable, and fixes every identity
which is incident with $T$ in the sense of \noderef{TreeIncidentEdge}.

We next justify a fact about the defective color which will be used in the
path argument.  Membership of $d$ in \noderef{DefectiveColors}, together
with \noderef{DefectiveColor}, gives at least two, and hence at least one,
identity $b$ in the $d$-colored boundary set.  By
\noderef{BoundaryColorEdges}, $b$ is a boundary identity of $V(T)$ and
$\varphi(b)=d$.  If $d$ were missing at any $x\in V(T)$, applying
closedness as defined in \noderef{ColoredVertexSetClosed} to $b$ and $x$
would be a contradiction.  Thus $d$ is not missing at any tree vertex.
Equivalently, by \noderef{MissingColors} and \noderef{PresentColors}, at
every $x\in V(T)$ there is an identity incident with $x$ which has color
$d$.  In particular $d$ is present at $y_0$ and at $v_0$.  Since
$\alpha_0$ is missing at $y_0$ and $\gamma$ is missing at $v_0$, we have
$\alpha_0\ne d$ and $\gamma\ne d$.

Let $P$ be the recorded $(\alpha_0,d)$ path, let $z$ be its terminal
vertex, and let $R$ be the set of its nodes at indices at least the
recorded index of $v_0$.  The full-path field of
\noderef{ScheideDefectiveSetup}, interpreted using
\noderef{ScheideDefectiveChainPath}, says that $P$ is simple, that its
support is the entire old $(\alpha_0,d)$-component of $y_0$, and that $z$
is outside $V(T)$.  The recorded successive indices identify the first
three nodes of $R$ as $v_0,f^0_d,u^0_d$.  Its last-tree-vertex field says
that every later vertex-node is outside $V(T)$.  Consequently the
vertex-nodes of $R$ meet $V(T)$ in exactly $v_0$.

Every edge-node of $P$ has old color $\alpha_0$ or $d$ by the adjacency
definition \noderef{TwoColorIncidenceGraph}.  At any vertex there is at
most one incident identity of either fixed color: two distinct such
identities would be adjacent by \noderef{AdjacentEdges} and would
contradict \noderef{IsProperEdgeColoring}.  Since
$\alpha_0\ne d$, the two path identities at each internal vertex of $P$
therefore have opposite colors.  At $y_0$, the color $\alpha_0$ is
missing, whereas $d$ is present by the preceding paragraph.  Thus the
first identity of $P$ is the unique incident $d$-identity, the old
two-color degree of $y_0$ is one, and the colors on $P$ alternate between
$d$ and $\alpha_0$.

The first edge of $R$ is $f^0_d$, which has color $d$ by
\noderef{ScheideDefectiveSetup}; it is incident with $v_0\in V(T)$ and
is fixed.  Every subsequent edge-node of $R$ lies between two later
vertex-nodes of $P$, both outside $V(T)$, and hence is not incident with
$T$.  The exact edgewise definition \noderef{OutsideTreeColorSwap} then
changes its old value $\alpha_0$ to $\gamma$ and leaves its old value $d$
unchanged, since $d$ is distinct from both swapped colors.  It follows
that $R$ is a $(\gamma,d)$ walk for $\varphi_1$.

We show that this walk has no new two-color adjacency at its vertices.
At $v_0$, every incident identity is fixed and $\gamma$ is missing.
Properness makes $f^0_d$ the unique incident $d$-identity.  If $v_0$ has
a predecessor edge on $P$, alternation makes its old color $\alpha_0$;
that edge is incident with $v_0\in V(T)$, so it remains $\alpha_0$ and is
not a $(\gamma,d)$-edge.  Hence $v_0$ has new two-color degree one.  At
each internal vertex of $R$, the two path identities have the two
distinct colors $\gamma,d$.  The already established properness of
$\varphi_1$, together with \noderef{AdjacentEdges}, excludes every third
incident identity of either color.

It remains to check the terminal vertex $z$.  Apply
\noderef{TwoColorIncidenceDegreeAtMostTwo} to the old proper coloring:
every node of the old two-color graph has degree at most two.  The finite
connected component is the support of the simple path $P$, and its
initial node $y_0$ has degree one.  Its terminal node $z$ must also have
degree one.  Indeed, an additional adjacency at $z$ cannot lead outside
the support because the support is the whole component.  Its other endpoint
would therefore be an edge-node already on $P$.  Every edge-node on this
vertex-to-vertex path is internal and already has its two distinct path
neighbors.  Unless it is the final path edge (which gives the adjacency
already counted at $z$), the extra adjacency would give that edge-node
degree at least three, contrary to the degree bound.  Thus no additional
adjacency exists.

We now compare incident identities at $z$ before and after the swap.
Because $z\notin V(T)$, any identity incident with both $z$ and $T$ is a
boundary identity by \noderef{BoundaryEdges} and
\noderef{TreeIncidentEdge}.  Closedness
\noderef{ColoredVertexSetClosed}, using the tree witnesses at which
$\alpha_0$ and $\gamma$ are missing, forbids either of those colors on
such a boundary identity.  Therefore every old $\alpha_0$-identity at
$z$ is not incident with $T$ and is changed to $\gamma$ by the exact
formula in \noderef{OutsideTreeColorSwap}; every old $d$-identity is
fixed because $d\ne\alpha_0,\gamma$.  Conversely, a new $\gamma$-identity
at $z$ cannot be tree-incident (a tree-incident identity is fixed, and an
old boundary $\gamma$-identity is impossible), so the same formula says
that it was formerly an $\alpha_0$-identity.  A new $d$-identity was
formerly a $d$-identity.  Thus the new $\gamma/d$ incidences at $z$ are
exactly the old $\alpha_0/d$ incidences, identity by identity.  The new
two-color degree of $z$ is consequently one as well.

Finally, every edge-node of $R$ is internal to this vertex-to-vertex
path.  Its two path neighbors are its two endpoint vertices, and the
loopless endpoint pair in \noderef{LooplessMultigraph} admits no further
incidence neighbor.  The vertex analysis above therefore shows that the
support of $R$ is connected and closed under every adjacency in the new
$(\gamma,d)$ incidence graph.
Since it contains $v_0$, the definition \noderef{TwoColorComponent}
gives that its support is exactly $K$.  This is precisely the asserted
suffix-set equality; it also puts $f^0_d$ in $K$ and proves that the only
tree vertex-node in $K$ is $v_0$.

If a displayed edge of $T$ belonged to $R$, then its two path-neighbor
vertices would be its two distinct endpoints.  Both endpoints of a
displayed edge lie in $V(T)$ by \noderef{TreeSequence} and
\noderef{TreeSequenceVertexSet}, contradicting the fact that $R$ has only
one tree vertex.  Hence no displayed edge identity belongs to $K$.
Lastly, any edge-node in this nonisolated component occurs next to a
vertex-node along a component path; by \noderef{TwoColorIncidenceGraph}
its $\varphi_1$-color is $\gamma$ or $d$.  All claimed assertions follow.
\end{proof}
